\chapter{Deep Dive: The Gateway Under Load --- Reactive Internals, Rate Limits, and the Edge's Job}
\label{ch:dd-gateway}

Chapter~4 taught the gateway as a route table: predicates, \code{lb://}
URIs, \code{StripPrefix}, one JWT filter. That is the gateway on a quiet
day. This chapter is the gateway on a bad one --- a slow course-service,
a client in a retry loop, a pod restart at the wrong moment --- and it
asks what the reactive stack actually buys, what the project's edge is
missing, and where the line runs between an edge's job and everybody
else's.

Three questions organize the chapter:

\begin{enumerate}
  \item \textbf{Why is the gateway non-blocking}, and what exactly goes
  wrong if a filter blocks?
  \item \textbf{What breaks first under load} --- and why the answer is
  ``the thing without a timeout''?
  \item \textbf{What belongs at the edge} --- rate limits, circuit
  breakers, CORS --- and what never should?
\end{enumerate}

\section{Why a gateway, and what it solves}

Without the gateway the frontend would hold six addresses, each service
would verify JWTs, and CORS, rate limits and tracing headers would be
implemented six times or zero times. The gateway collapses that into one
address and one place for every \emph{cross-cutting} concern --- concerns
that are about \emph{every} request rather than about courses or users.
The interview form of the question is ``what does an API gateway do that
a load balancer does not?'' and the senior answer is: a load balancer
picks a backend; a gateway \emph{understands the request} --- it can
authenticate, rewrite, rate-limit, and fail over per route, because it
runs application code on the path. That power is also the risk: whatever
you put on that path, every request pays for.

\begin{decision}[an application gateway behind an infrastructure router]
The cluster already has Traefik (Chapter~18). Why keep Spring Cloud
Gateway behind it? Because the two do different jobs: Traefik owns the
host, TLS, and the path fan-out; the gateway owns things only the
application knows --- which Eureka name serves \code{/api/courses}, what
a valid JWT looks like, which headers to inject. Merging them would mean
teaching Traefik the JWT secret and the header contract, which is
possible (ForwardAuth middleware) but moves application security into
infrastructure config. The decision flips if the platform provides
identity itself (a cloud API Gateway with Cognito or IAP): then the
Spring gateway is a second hop paid for nothing.
\end{decision}

\section{The reactive stack, read from the filter}
\label{sec:reactive}

The gateway declares \code{web-application-type: reactive} in its
\code{application.yml}. That one line swaps Tomcat's thread-per-request
model for Reactor Netty's event loop, and every line of the JWT filter
is written for that model:

\begin{lstlisting}[language=Java]
@Override
public Mono<Void> filter(ServerWebExchange exchange,
                         GatewayFilterChain chain) {           (*@\co{1}@*)
    String path = exchange.getRequest().getURI().getPath();
    if (isPublicPath(path)) {
        return chain.filter(exchange);                          (*@\co{2}@*)
    }
    String authHeader = exchange.getRequest().getHeaders()
            .getFirst(HttpHeaders.AUTHORIZATION);
    if (authHeader == null || !authHeader.startsWith("Bearer ")) {
        return unauthorizedResponse(exchange, "Missing ...");  (*@\co{3}@*)
    }
    String token = authHeader.substring(7);
    try {
        Claims claims = Jwts.parser().verifyWith(signingKey)
                .build().parseSignedClaims(token).getPayload(); (*@\co{4}@*)
        ServerHttpRequest mutated = exchange.getRequest().mutate()
                .header(JwtConstants.HEADER_USER_ID, claims.getSubject())
                .header(JwtConstants.HEADER_USER_ROLE,
                        claims.get(JwtConstants.CLAIM_ROLE, String.class))
                .build();
        return chain.filter(exchange.mutate().request(mutated).build());
    } catch (Exception e) {
        return unauthorizedResponse(exchange, "Invalid or expired token");
    }
}
\end{lstlisting}

\begin{description}
  \item[\co{1}] The return type is the whole story. A servlet filter
  returns \code{void} \emph{after} the request is done; this returns a
  \code{Mono<Void>} --- a \emph{description} of work that completes
  later. The thread that called \code{filter} is free the moment this
  method returns.
  \item[\co{2}] \code{chain.filter(exchange)} does not forward the
  request; it returns the next stage's \code{Mono}. Nothing has happened
  yet. The chain is assembled first and subscribed once.
  \item[\co{3}] Even the error path is reactive: \code{writeWith} returns
  a \code{Mono} that completes when Netty has flushed the body.
  \item[\co{4}] The one piece of real work, and it is CPU-bound and
  microseconds long: an HMAC over a few hundred bytes. This is the kind
  of work that \emph{may} run on an event-loop thread. A database call
  is not.
\end{description}

\subsection{Threads: four, not four hundred}

Tomcat answers concurrency with threads: 200 by default, one per
in-flight request, each parked while its backend answers. A proxy's
requests are almost entirely waiting, so a Tomcat proxy under 1{,}000
in-flight requests needs 1{,}000 threads and about a gigabyte of stacks.
Reactor Netty runs a small fixed pool of \emph{event-loop} threads, each
multiplexing thousands of sockets with non-blocking I/O: a request
arrives, the filter chain is assembled, the upstream write is issued, and
the thread moves on. When the upstream answers, whichever loop owns that
socket resumes the chain.

\begin{hardware}
Reactor Netty sizes the loop as \code{max(availableProcessors, 4)}.
On ts-server's two cores that is \textbf{four} event-loop threads for
the whole gateway. Four threads serving every request is the reason the
gateway's Deployment asks for less memory than an MVC service
(384\,Mi vs.\ 512\,Mi, Chapter~17) --- and the reason one blocked thread
is not a 0.5\,\% problem but a 25\,\% one.
\end{hardware}

\subsection{What blocking does to an event loop}

Suppose someone ``improves'' the filter by checking a deny-list in
Postgres:

\begin{lstlisting}[language=Java]
// DO NOT: a JDBC call inside a WebFlux filter
boolean revoked = jdbcTemplate.queryForObject(
        "select exists(select 1 from revoked where jti = ?)",
        Boolean.class, claims.getId());
\end{lstlisting}

JDBC is blocking by specification: the calling thread waits for the
socket. That thread is one of four event loops, and while it waits it
is not servicing the hundreds of other sockets it owns. Their requests
stall --- not fail, \emph{stall} --- until the query returns. Under a
slow database, all four loops are pinned within seconds and the gateway
accepts connections it cannot process: every route, every user, one
symptom. The same is true of \code{.block()} on a \code{Mono},
\code{Thread.sleep}, a synchronous \code{RestTemplate} call, or a file
read. The reactive alternative moves the work off the loop and
\emph{composes} it into the chain:

\begin{lstlisting}[language=Java]
return Mono.fromCallable(() -> deny.isRevoked(claims.getId()))
        .subscribeOn(Schedulers.boundedElastic())   // off the loop
        .flatMap(revoked -> revoked
                ? unauthorizedResponse(exchange, "Token revoked")
                : chain.filter(exchange.mutate().request(mutated).build()));
\end{lstlisting}

Better still is a non-blocking client (reactive Redis, R2DBC) so no
thread waits at all. The rule for anything on the gateway path:
\emph{if the library's method returns a value instead of a}
\code{Mono}\emph{/}\code{Flux}\emph{, it does not belong in a filter.}
Reactor's BlockHound agent turns that rule into a test failure.

\subsection{Predicates, order, and the \code{StripPrefix} arithmetic}

Routes are tried in declaration order (all have the default
\code{order} of 0, and the sort is stable); the first predicate match
wins. The project's table has no overlaps, but a future
\code{Path=/api/**} catch-all placed \emph{above} the specific routes
would swallow them silently --- put catch-alls last or give them a high
\code{order}. Filters attached to a route are numbered in sequence and
run \emph{after} global filters with a lower order; the JWT filter's
\code{getOrder()} of $-1$ places it before every route filter, which is
what lets a route's rate limiter (below) read the \code{X-User-Id}
header the JWT filter just wrote.

\begin{pitfall}[the doubled segment]
Symptom: the frontend calls \code{/api/courses/courses/5} and it
works, but a hand-written \code{/api/courses/5} returns 404 from
course-service. Cause: \code{StripPrefix=2} removes exactly two
segments, \code{/api} and \code{/courses}, so \code{/api/courses/5}
arrives as \code{/5} --- and no controller is mapped there. The first
\code{courses} selects the \emph{route}; the second survives the strip
and matches \code{@RequestMapping("/courses")}. The frontend repeats the
segment on purpose. The cleaner design is \code{StripPrefix=1} with
controllers keeping their resource prefix; the project chose to keep
services unaware of \code{/api} at the cost of the doubled URL.
\end{pitfall}

\section{CORS at the edge}

\code{CorsConfig} registers one \code{CorsWebFilter} for \code{/**}:
origins \code{localhost:3000} and \code{5173}, six methods, any header,
credentials allowed. Three things to understand about where it lives.

\emph{Why here and not in each service.} CORS is a browser rule about
\emph{origins}, and the origin the browser talks to is the gateway.
A service answering \code{Access-Control-Allow-Origin} would be
answering a question nobody asked it; and if two services disagreed, the
browser would see inconsistent policy for one API. Edge concerns go at
the edge.

\emph{What a preflight costs.} For any request that is not ``simple''
--- a JSON \code{POST}, anything with an \code{Authorization} header ---
the browser first sends an \code{OPTIONS} request and waits for the
answer. That is one extra round trip per call, and every API call in
this project carries \code{Authorization}. The gateway answers preflights
itself without touching a route, which is cheap, but the round trip is
paid by the user. \code{Access-Control-Max-Age} lets the browser cache
the answer; \code{CorsConfig} does not set it, so Chrome falls back to
five seconds. One line, \code{config.setMaxAge(3600L)}, removes almost
every preflight from a session.

\emph{When it is never consulted.} In the cluster the frontend and
\code{/api} share one origin behind Traefik (Chapter~26), so no request
is cross-origin and the filter is idle. Dev is the only place it runs
--- and Vite's proxy could make even dev same-origin. CORS config is
therefore not a security feature here; it is a development convenience
that must not be widened carelessly: \code{allowCredentials(true)} with
a wildcard origin is refused by browsers precisely because it would let
any site call the API with the user's cookies.

\section{The edge's job, and what is not}

\begin{center}
\small
\begin{tabular}{@{}p{0.44\linewidth}p{0.48\linewidth}@{}}
\toprule
Belongs at the edge & Does not \\
\midrule
Authentication, identity headers & Authorization rules per resource \\
Rate limiting, quotas & Business validation \\
Timeouts, retries (idempotent), circuit breaking & Response aggregation across services \\
CORS, TLS termination, compression & Calling a database \\
Tracing headers, access logs & Data transformation beyond headers/paths \\
\bottomrule
\end{tabular}
\end{center}

The right-hand column is the ``smart pipes'' anti-pattern the ESB era
taught: once the gateway knows business rules, every feature ships in
two places and the gateway team becomes the bottleneck for everyone. The
one legitimate exception --- composing several service calls into one
screen-shaped response --- is a service of its own (a \emph{backend for
frontend}), deployed and scaled separately from the edge.

\section{Pros and cons, honestly}

\begin{itemize}
  \item \textbf{Pro:} one address, one JWT filter, one CORS policy;
  routes are configuration, changeable without a rebuild
  (Chapter~2).
  \item \textbf{Pro:} the reactive model holds thousands of in-flight
  requests on four threads and boots in seconds.
  \item \textbf{Con:} the project's gateway sets \emph{no} HTTP client
  timeouts. Reactor Netty's defaults are a 30-second connect timeout and
  \emph{no} response timeout at all.
  \item \textbf{Con:} no circuit breaker, no rate limiter --- every
  client can send as fast as it likes, and a dead service is discovered
  by each request individually.
  \item \textbf{Con:} \code{replicas: 1} in Kubernetes. The component
  every request passes is the one with a single copy.
  \item \textbf{Con:} the starter it uses, \code{spring-cloud-starter-gateway},
  is deprecated in Spring Cloud 2025.0 in favor of
  \code{spring-cloud-starter-gateway-server-webflux}; the rename is
  harmless today and a build break in a future upgrade.
\end{itemize}

\section{What breaks}

\begin{pitfall}[a slow downstream takes the whole route with it]
Symptom: course-service is slow (a lock, a full GC), and within a minute
\emph{all} requests to \code{/api/courses/**} fail with 503 or hang ---
including cheap ones that would have succeeded. Cause: the gateway keeps
a per-host connection pool, sized \code{max(cores, 8) $\times$ 2} = 16
on this box. With no response timeout, each slow request holds one
pooled connection for as long as course-service takes; the seventeenth
request waits in the acquire queue for 45 seconds and then fails. Every
user is queued behind the slowest one. Fix: a response timeout
(Section~\ref{sec:gw-enhance}) so a stuck connection is reclaimed, and a
circuit breaker so the route fails fast once failures cluster.
\end{pitfall}

\begin{pitfall}[431 Request Header Fields Too Large]
Symptom: a user with many cookies, or a token with fat claims
(Chapter~28), gets an immediate 431 from the gateway with nothing in the
service logs. Cause: Netty caps the request \emph{header block} at
8\,KB by default; the gateway rejects the request before any filter
runs. Fix: keep tokens thin. Raising
\code{server.max-http-request-header-size} is the second-best answer,
because every hop behind the gateway has its own cap.
\end{pitfall}

\begin{pitfall}[the gateway is the single point of failure]
Symptom: a rolling update of the gateway, a node reboot, or an OOM kill
--- and the entire application is down, while every service behind it
is healthy. Cause: one replica. The gateway is stateless, so the fix is
\code{replicas: 2} plus a \code{PodDisruptionBudget}; the only thing
that would make scaling harder is per-instance state such as an
in-memory rate limiter, which is exactly why the enhancement below keeps
that state in Redis.
\end{pitfall}

\section{What breaks at 3\,a.m.}

course-service's Postgres starts a long autovacuum; queries that took
20\,ms take 4\,s. Walk the path.

\begin{enumerate}
  \item The frontend fires a dashboard's six requests. Four hit
  course-service and take 4\,s each. The gateway holds four pooled
  connections for 4\,s.
  \item Users retry. Twenty users, four requests each: 80 requests
  compete for 16 connections. Acquire queue grows; after 45\,s waiting
  requests fail with 503.
  \item Requests to auth-service and dictionary-service are
  \emph{unaffected} --- the pool is per host --- until users' access
  tokens expire and \code{/refresh} succeeds but the dashboard still
  fails, so they log out and back in, adding load.
  \item Nobody is paged, because the gateway is ``healthy'':
  \code{/actuator/health} answers 200 and its event loops are not
  blocked. The failure is invisible at the edge and visible only as
  latency in Zipkin (Chapter~24).
\end{enumerate}

With a 5\,s response timeout and a circuit breaker on the route, step~2
becomes: slow calls are cut at 5\,s and counted as failures; after ten
of them the breaker opens; for the next 10\,s every \code{/api/courses}
request fails in under a millisecond with a fallback body; the pool is
never exhausted; auth and dictionary keep working; the breaker's state
change is a metric you can alert on. The database is still slow ---
resilience does not fix causes --- but the blast radius is one route
instead of the site.

\section{Enhancing the resolution}
\label{sec:gw-enhance}

Three additions, all configuration plus one bean. Dependencies first
(gateway \code{pom.xml}):

\begin{lstlisting}[language=yaml]
# spring-cloud-starter-circuitbreaker-reactor-resilience4j
#   -> the reactive CircuitBreaker gateway filter
# spring-boot-starter-data-redis-reactive
#   -> backing store for RequestRateLimiter (non-blocking client)
\end{lstlisting}

Then \code{api-gateway.yml} in the config-server:

\begin{lstlisting}[language=yaml]
spring:
  data:
    redis:
      host: ${REDIS_HOST:localhost}            # redis:6379 in compose
  cloud:
    gateway:
      httpclient:
        connect-timeout: 2000                  (*@\co{1}@*)
        response-timeout: 5s
        pool:
          max-connections: 64                  (*@\co{2}@*)
          acquire-timeout: 2000
      default-filters:
        - RemoveRequestHeader=X-User-Id        (*@\co{3}@*)
        - RemoveRequestHeader=X-User-Role
      routes:
        - id: course-service-courses
          uri: lb://course-service
          predicates:
            - Path=/api/courses/**
          filters:
            - StripPrefix=2
            - name: CircuitBreaker             (*@\co{4}@*)
              args:
                name: courseCb
                fallbackUri: forward:/fallback/course
                statusCodes: "500,502,503,504"
            - name: RequestRateLimiter         (*@\co{5}@*)
              args:
                redis-rate-limiter.replenishRate: 20
                redis-rate-limiter.burstCapacity: 40
                redis-rate-limiter.requestedTokens: 1
                key-resolver: "#{@userKeyResolver}"
resilience4j:
  circuitbreaker:
    instances:
      courseCb:
        sliding-window-size: 10
        minimum-number-of-calls: 10
        failure-rate-threshold: 50
        slow-call-duration-threshold: 3s       (*@\co{6}@*)
        slow-call-rate-threshold: 80
        wait-duration-in-open-state: 10s
        permitted-number-of-calls-in-half-open-state: 3
  timelimiter:
    instances:
      courseCb:
        timeout-duration: 5s                   (*@\co{7}@*)
\end{lstlisting}

\begin{description}
  \item[\co{1}] Connect in 2\,s or give up; a full response within 5\,s
  or the connection is closed and returned to the pool. Per-route
  overrides are possible via route \code{metadata} for a known-slow
  endpoint such as file upload.
  \item[\co{2}] Sixteen connections per host is tight once timeouts
  make requests short; 64 is comfortable on this box and still bounded.
  \item[\co{3}] The Chapter~28 fix: identity headers can only be
  written by the gateway.
  \item[\co{4}] The reactive Resilience4j filter. Failures (and the
  listed status codes) count toward the breaker; when open, requests are
  forwarded to \code{/fallback/course} \emph{inside the gateway}.
  \item[\co{5}] Token bucket in Redis: 20 tokens per second refill, up
  to 40 stored, one spent per request. Exceeding it returns 429 with
  \code{X-RateLimit-Remaining: 0}. Redis, not memory, so two gateway
  replicas share one bucket per user.
  \item[\co{6}] Without this, slow calls are not failures --- the
  default slow-call threshold is 60\,s. Three seconds means ``slower
  than we designed for''.
  \item[\co{7}] The gateway's breaker wraps every call in a
  TimeLimiter whose default is \textbf{one second}. Left unset, a
  healthy 1.2\,s upload would trip the breaker. Set it explicitly and
  keep it $\ge$ the response timeout.
\end{description}

The key resolver and the fallback endpoint:

\begin{lstlisting}[language=Java]
@Configuration
public class RateLimitConfig {

    /** Bucket per authenticated user; anonymous callers share
     *  a bucket per client IP (public paths: login, refresh). */
    @Bean
    public KeyResolver userKeyResolver() {
        return exchange -> {
            String user = exchange.getRequest().getHeaders()
                    .getFirst(JwtConstants.HEADER_USER_ID);
            if (user != null) {
                return Mono.just("u:" + user);
            }
            var addr = exchange.getRequest().getRemoteAddress();
            return Mono.just("ip:" + (addr == null ? "unknown"
                    : addr.getAddress().getHostAddress()));
        };
    }
}

@RestController
public class FallbackController {

    /** Served by the gateway itself when a breaker is open.
     *  503 + Retry-After tells clients (and browsers) to back off. */
    @GetMapping("/fallback/course")
    public ResponseEntity<Map<String, String>> course() {
        return ResponseEntity.status(HttpStatus.SERVICE_UNAVAILABLE)
                .header(HttpHeaders.RETRY_AFTER, "10")
                .body(Map.of("error", "course-service unavailable",
                             "retryAfterSeconds", "10"));
    }
}
\end{lstlisting}

\begin{decision}[what to key the rate limit on]
Per \emph{user id} is the right default for an authenticated API: it
matches the thing you are protecting (a tenant's fair share) and cannot
be dodged by rotating IPs. Per \emph{IP} is the only option before
authentication --- login and refresh --- and is the right one there,
because the threat is credential stuffing from one source; it is wrong
behind a corporate NAT, where one IP is a thousand users. Per \emph{API
key} is for machine clients and partners, where the key is the contract.
Production edges usually stack two: a coarse per-IP limit on public
paths and a per-user limit on everything else.
\end{decision}

\begin{handson}[hit the bucket]
With Redis added to \code{docker-compose.yml} and the gateway restarted,
fire 50 requests as one user faster than the bucket refills:
\begin{lstlisting}[language=bash]
$ for i in $(seq 1 50); do
    curl -s -o /dev/null -w "%{http_code} " \
      -H "Authorization: Bearer $TOKEN" \
      http://localhost:8080/api/courses/courses
  done; echo
200 200 200 ... 200 429 429 429 429 429 429 429 429 429 429
$ curl -si -H "Authorization: Bearer $TOKEN" \
    http://localhost:8080/api/courses/courses | grep -i ratelimit
X-RateLimit-Remaining: 0
X-RateLimit-Burst-Capacity: 40
X-RateLimit-Replenish-Rate: 20
\end{lstlisting}
Forty succeed (the burst), then 429s until a second passes and twenty
tokens return. Stop course-service and repeat: after ten failures the
same command returns the fallback JSON in under a millisecond, and
\code{/actuator/circuitbreakers} shows \code{courseCb} as \code{OPEN}.
\end{handson}

\section{Outside the project}

Kong, Envoy (and Istio's ingress built on it), nginx, and the cloud
gateways (AWS API Gateway, Apigee) implement the same four ideas ---
timeouts, breakers, rate limits, header rewriting --- as native
configuration rather than Spring filters, with the token bucket
typically in the proxy's own memory per node plus a shared store for
global limits. Envoy's ``outlier detection'' is the circuit breaker;
its ``local rate limit'' filter is the bucket. What you have learned
here transfers directly: the questions an interviewer asks about any of
them --- where does state live, what happens when the store is down,
what is the timeout hierarchy --- have the same answers.

\section{Questions from real interviews}

\subsection*{Q: Traffic goes 10$\times$ on Monday morning. What breaks first in the gateway, and what do you tune?}

Not the event loops. Four loops doing non-blocking I/O handle far more
than ten times this project's traffic; CPU on the gateway is HMAC
verification and header copying, which is microseconds. What breaks
first is whatever is \emph{unbounded}: the per-host connection pool (16
by default here) fills as soon as any downstream slows, and with no
response timeout it never drains. So the tuning order is: response and
connect timeouts so connections are reclaimed; the pool ceiling raised
to something the downstream can actually serve (\code{max-connections},
and a short \code{acquire-timeout} so queued requests fail fast rather
than pile up); a circuit breaker per route so a struggling service is
shed instead of queued; then horizontal scaling --- two or three
replicas behind the Service with an HPA on CPU, which for a stateless
gateway is trivial once rate-limit state is in Redis. Memory limits
matter too: Netty keeps direct buffers off-heap, so a 512\,Mi limit
with \code{MaxRAMPercentage=75} leaves 128\,Mi for them; under 10$\times$
buffering that is where an OOM kill would come from.

\subsection*{Q: One downstream is slow and the whole site hangs. Why, and what do you do?}

Because the site shares the thing that is exhausted. If the gateway has
one connection pool per host, only that host's route should degrade ---
so ``the whole site'' says something else is shared: an event loop
blocked by a synchronous call in a filter, a global rate limiter keyed
on the wrong thing, or the frontend firing all its calls in one
\code{Promise.all} so one slow response holds the page. The fix is in
that order: audit filters for blocking calls (BlockHound in tests),
give the slow route a timeout and breaker so it fails in milliseconds
with a fallback body, and make the frontend render partial pages. The
senior point is that resilience is about \emph{isolation}: the goal is
not that course-service never slows down but that its slowness cannot
reach the login page.

\subsection*{Q: How do you rate-limit per user across three gateway replicas?}

State must be shared or the limit is per replica and off by a factor
of three. Spring Cloud Gateway's \code{RequestRateLimiter} runs a token
bucket as a Lua script in Redis, atomically, keyed by whatever the
\code{KeyResolver} returns --- here \code{u:<userId>} from the header
the JWT filter set, falling back to the client IP on public paths. The
numbers are the policy: \code{replenishRate} is sustained requests per
second, \code{burstCapacity} what a page load may spend at once; set
burst to roughly one screen's worth of calls. Redis becomes a
dependency, so decide its failure mode deliberately: the filter's
default is to \emph{allow} when Redis is unreachable, which is right
for an internal app and wrong for a public one facing abuse. For very
high volume, a local per-replica bucket in front of the shared one
absorbs most of the Redis round trips.

\subsection*{Q: Where do you put authentication, and where do you not put logic?}

Authentication --- ``is this token valid and whose is it'' --- at the
edge, once, because it is the same question for every request and
because that is where the credential enters the system. Authorization
--- ``may this user grade this submission'' --- in the service that owns
submissions, because only it knows the data. The gateway may enforce
coarse rules that are still request-shaped (``\code{/api/admin/**}
requires role ADMIN'') but the moment a rule needs a database row it
belongs downstream. Nothing business-shaped goes in the gateway:
no aggregation, no transformation of bodies, no calls to data stores.
The test I use in reviews: if a change to a domain rule requires a
gateway deploy, the line has been crossed.

\subsection*{Q: How do you deploy the gateway with zero dropped requests?}

Two replicas at minimum, a \code{PodDisruptionBudget} of
\code{minAvailable: 1}, and a rolling update with
\code{maxUnavailable: 0} so the new pod is Ready before an old one is
removed. Readiness must be real: the probe should fail while Eureka's
registry is still empty, or the new pod is routed to before it can
route anywhere. On shutdown, the old pod must stop receiving \emph{and}
finish in-flight requests: \code{server.shutdown: graceful} with a
\code{spring.lifecycle.timeout-per-shutdown-phase} longer than the
response timeout, and a \code{preStop} sleep of a few seconds so the
kube-proxy endpoint removal propagates before the process gets
SIGTERM. That last gap --- the endpoint list and the pod's lifetime are
updated by different controllers --- is what drops the requests people
blame on ``the rolling update''.

\subsection*{Q: A partner integration needs one response built from three services. Gateway?}

No. That is a backend-for-frontend: a small service that calls the three
in parallel, tolerates one failing, and shapes the response for that
client. It scales and deploys on its own schedule, owns its own
timeouts, and keeps the gateway a proxy. Put it in the gateway and every
partner change is an edge deploy, and the gateway's event loops start
holding fan-out state for the slowest of three calls.

\begin{quiz}
\begin{enumerate}
  \item The JWT filter returns \code{chain.filter(exchange)} without
  awaiting anything. Explain what has and has not happened at the moment
  that line returns, and why a servlet filter could not be written this
  way.
  \item A colleague adds a synchronous Redis lookup to the JWT filter to
  check a deny-list. Under a 200\,ms Redis stall, what happens to
  requests for \emph{other} routes, and why is the impact on this
  machine 25\,\% per stalled call?
  \item course-service hangs on 10\,\% of requests. With the project's
  current gateway configuration, how many concurrent hung requests
  block \emph{all} traffic to that service, and which single setting
  changes that number?
  \item The gateway's \code{CircuitBreaker} filter is added to a route
  without a \code{timelimiter} block. A healthy upload takes 1.5\,s.
  What happens, and why?
  \item Why is the rate limiter's state in Redis rather than in the
  gateway's memory, and what is the trade-off when Redis is down?
  \item Every request from the frontend carries \code{Authorization}.
  What does that imply for CORS preflights in development, what single
  setting reduces them, and why does none of this apply in the cluster?
\end{enumerate}
\end{quiz}
